% A l'original la figura és a la dreta de l'apartat a.
\begin{solapartat}
\figura[0.3]{sol_fig1.png}

\punts{0,4} $\dfrac{F_e}{F_g} = \dfrac{eE}{mg}$

\medskip
% La pauta fa servir g = 9,80 (l'enunciat dona 9,81 m s^-2); el resultat no canvia.
\punts{0,4} $\dfrac{F_e}{F_g} = \dfrac{1,60\times 10^{-19}\times 2000}{9,10\times 10^{-31}\times 9,80} = 3,6\times 10^{13}$

\medskip
\punts{0,2} La força elèctrica és $3,6\times 10^{13}$ vegades més gran que la força gravitatòria.
\end{solapartat}

\begin{solapartat}
L'electró trigarà un temps per a recórrer la distància de 1 cm en la direcció \textit{x}:

\medskip
\punts{0,4} $t = \dfrac{x}{v_0} = \dfrac{10^{-2}}{10^{6}} = 10^{-8}\,s$

\smallskip
En aquest temps, l'electró és desviat cap amunt, una distancia ”\textit{y}” antiparal·lela al camp elèctric:

\medskip
\punts{0,3} $F_e = eE = ma_y \Rightarrow a_y = \dfrac{eE}{m} = 3,52\times 10^{14}\,ms^{-2}$

\smallskip
\punts{0,3} $y = \dfrac{1}{2}at^2 = 1,76\times 10^{-2}\,m = 1,76\,cm$
\end{solapartat}
